\documentclass[conference]{IEEEtran}

\usepackage{graphicx}
\usepackage{amsmath,amssymb}
\usepackage{float}
\usepackage{booktabs}
\usepackage{multirow}

\begin{document}

\title{
    \textbf{Benchmark Report: 507.cactuBSSN\_r} \\
    \text{Register Spill Analysis on RISC-V via gem5}
}

\author{
    \IEEEauthorblockN{Lütfullah Burak Kaya}
    \IEEEauthorblockA{
        Ozyegin University \\
        burak.kaya.50711@ozu.edu.tr
    }
    \and
    \IEEEauthorblockN{Ismail Akturk}
    \IEEEauthorblockA{
        Ozyegin University \\
        ismail.akturk@ozyegin.edu.tr
    }
}

\newcommand{\LongDate}{%
    \number\day\space
    \ifcase\month
    \or January\or February\or March\or April\or May\or June%
    \or July\or August\or September\or October\or November\or December%
    \fi
    \space \number\year
}

\twocolumn[
\begin{@twocolumnfalse}
\maketitle
\begin{center}{\small \LongDate}\end{center}
\vspace{1em}
\end{@twocolumnfalse}
]

% ================================================

\begin{abstract}
This report presents the register spill characterization of the
\texttt{507.cactuBSSN\_r} benchmark from SPEC CPU2017, executed on a
RISC-V (RV64GC) target using the gem5 TimingSimpleCPU simulator.
A custom SpillDetector was integrated into gem5 to count store-then-load
patterns on the stack within a user-defined Region of Interest (ROI).
The test dataset (\texttt{spec\_test.par}, grid size $42^3$, 10 time
steps) was used as the input.
Results show a spill rate of approximately \textbf{7.88\%} within the
ROI --- the highest observed across all SPEC CPU2017 benchmarks
evaluated in this study --- corresponding to 3.78 billion detected
spill events across 48 billion ROI instructions.
The benchmark's stencil-based numerical relativity kernel creates
extreme register pressure due to the large number of simultaneously
live floating-point variables required to evaluate the BSSN equations
at each grid point.
\end{abstract}

% =========================================================
\section{Benchmark Overview}

\texttt{507.cactuBSSN\_r} is the SPEC CPU2017 rate variant of
CactusBSSN, a numerical relativity application built on the
Cactus computational framework. It solves the
Baumgarte--Shapiro--Shibata--Nakamura (BSSN) formulation of the
Einstein field equations on a 3-dimensional uniform Cartesian grid
using a finite-difference stencil method. The benchmark is one of
the most floating-point-intensive workloads in SPEC CPU2017,
executing dense stencil sweeps over a 3D array with a large number
of evolved tensor fields.

The Cactus framework organises computation into \emph{thorns}
(independent modules). For this benchmark the active thorns are:
\texttt{PUGH} (uniform grid driver), \texttt{ML\_BSSN} (the McLachlan
BSSN evolution kernel), \texttt{MoL} (method-of-lines time integrator),
\texttt{ADMBase}, \texttt{CartGrid3D}, \texttt{GaugeWave}, and
supporting I/O and reduction thorns.

\subsection{ROI Instrumentation}

The RISC-V binary (\texttt{cactuBSSN\_r.riscv}) was compiled with gem5
ROI markers inserted in \texttt{src/Cactus/main/flesh.cc}, wrapping
the \texttt{CCTK\_Evolve()} call --- the top-level time-evolution
loop of the Cactus framework. The surrounding \texttt{CCTK\_Initialise()}
and \texttt{CCTK\_Shutdown()} calls are \emph{excluded} from the ROI,
so only the core numerical computation is measured:

\begin{itemize}
    \item \texttt{m5\_work\_begin(0,0)} --- before \texttt{CCTK\_Evolve()}
    \item \texttt{m5\_work\_end(0,0)}   --- after  \texttt{CCTK\_Evolve()}
\end{itemize}

\texttt{CCTK\_Evolve()} drives all time steps of the simulation, calling
the ML\_BSSN kernel repeatedly to advance the solution by one time step
per iteration.

% =========================================================
\section{Simulation Setup}

\subsection{Input Datasets}

All three dataset sizes use the same grid topology ($42^3$ or $84^3$
uniform Cartesian grid) and differ only in grid resolution and
iteration count (Table~\ref{tab:inputs}).

\begin{table}[H]
\centering
\caption{cactuBSSN\_r Input Dataset Sizes}
\label{tab:inputs}
\begin{tabular}{lccc}
\toprule
\textbf{Dataset} & \textbf{Grid size} & \textbf{Iterations} & \textbf{Parameter file} \\
\midrule
test    & $42^3$ & 10 & spec\_test.par \\
train   & $42^3$ & 80 & spec\_train.par \\
refrate & $84^3$ & 80 & spec\_ref.par \\
\bottomrule
\end{tabular}
\end{table}

The \textbf{test} dataset was selected for this run: a $42^3$ grid
evolved for 10 time steps. The grid contains $42 \times 42 \times 42 =
74{,}088$ points; at each point the BSSN equations update 25 evolved
fields (the conformal metric $\tilde{\gamma}_{ij}$, extrinsic curvature
$\tilde{A}_{ij}$, conformal factor $W$, trace $K$, and the shift vector
components). The refrate dataset increases both the linear grid size by
a factor of 2 (8$\times$ more grid points) and the iteration count by
8$\times$, making it approximately $64\times$ more expensive than the
test dataset.

\subsection{gem5 Configuration}

\begin{itemize}
    \item \textbf{CPU model:}   TimingSimpleCPU
    \item \textbf{ISA:}         RISC-V RV64GC
    \item \textbf{Caches:}      enabled (default L1 I/D + L2)
    \item \textbf{Memory:}      4\,GB simulated physical RAM
    \item \textbf{Spill mode:}  summary only (\texttt{spill\_verbose=False})
\end{itemize}

The full gem5 invocation was:

\begin{verbatim}
./build/RISCV/gem5.opt
  --outdir=benchmarks/cactuBSSN/m5out_test
  configs/deprecated/example/se.py
  --cmd=benchmarks/cactuBSSN/cactuBSSN_r.riscv
  --options="benchmarks/cactuBSSN/spec_test.par"
  --cpu-type=TimingSimpleCPU
  --caches
  --mem-size=4GB
\end{verbatim}

% =========================================================
\section{Simulation Results}

The simulation completed successfully. Cactus printed per-iteration
output showing convergence of the BSSN fields across all 10 time
steps, followed by \texttt{"Done."} and a clean exit.

\subsection{Pre-ROI Context}

Before the ROI was entered, Cactus executed full framework
initialisation: parameter parsing, thorn activation, grid construction,
and initial data setup (\texttt{CCTK\_Initialise()}). This pre-ROI
phase includes setting up all 74,088 grid points with GaugeWave
initial data.
Table~\ref{tab:preroi} shows the global counters at the moment
\texttt{m5\_work\_begin} was triggered.

\begin{table}[H]
\centering
\caption{Global Counters at ROI Entry}
\label{tab:preroi}
\begin{tabular}{lr}
\toprule
\textbf{Metric} & \textbf{Value} \\
\midrule
Total instructions (pre-ROI) & 936,822,146 \\
Total spills detected (pre-ROI) & 110,850,678 \\
Pre-ROI spill rate & 11.83\% \\
\bottomrule
\end{tabular}
\end{table}

The pre-ROI spill rate of 11.83\% is remarkably high, reflecting the
register pressure already present during Cactus framework initialisation
and initial data computation.

\subsection{ROI Spill Statistics}

Table~\ref{tab:roi} presents the spill statistics collected exclusively
within the ROI (10 BSSN time steps via \texttt{CCTK\_Evolve()}).

\begin{table}[H]
\centering
\caption{ROI Spill Statistics --- 507.cactuBSSN\_r (test, $42^3$, 10 iter.)}
\label{tab:roi}
\begin{tabular}{lr}
\toprule
\textbf{Metric} & \textbf{Value} \\
\midrule
ROI instructions         & 48,023,254,961 \\
ROI stores               &  6,919,103,141 \\
ROI loads                & 12,477,164,927 \\
\midrule
\textbf{ROI spills}      & \textbf{3,782,099,438} \\
\textbf{Spill rate}      & \textbf{7.8756\%} \\
\bottomrule
\end{tabular}
\end{table}

\subsection{Timing}

\begin{table}[H]
\centering
\caption{Simulation Timing}
\label{tab:timing}
\begin{tabular}{lr}
\toprule
\textbf{Metric} & \textbf{Value} \\
\midrule
Simulated time           & 263.59\,s \\
Host wall time           & 43,555\,s ($\approx$12.1\,hours) \\
Total simulated instructions & 48,961,298,333 \\
\bottomrule
\end{tabular}
\end{table}

% =========================================================
\section{Analysis}

\subsection{Spill Rate Interpretation}

The measured spill rate of \textbf{7.88\%} means that approximately
1 in every 13 instructions inside the ROI is a register spill.
This is the highest spill rate observed among all SPEC CPU2017
benchmarks evaluated in this study (Table~\ref{tab:compare}).

\begin{table}[H]
\centering
\caption{Cross-Benchmark Spill Rate Comparison}
\label{tab:compare}
\begin{tabular}{lrr}
\toprule
\textbf{Benchmark} & \textbf{ROI Spills} & \textbf{Spill Rate} \\
\midrule
507.cactuBSSN\_r & 3,782,099,438 & \textbf{7.8756\%} \\
500.perlbench\_r &     1,070,055 & 7.2838\% \\
526.blender\_r   &    61,845,869 & 6.2193\% \\
544.nab\_r       &   151,310,616 & 2.2612\% \\
505.mcf\_r       &   446,113,136 & 1.8263\% \\
508.namd\_r      &   327,718,955 & 1.0269\% \\
\bottomrule
\end{tabular}
\end{table}

\subsection{Store vs.\ Load Balance}

The load count (12.48B) is approximately $1.80\times$ higher than the
store count (6.92B). This ratio is consistent with a stencil-heavy
workload: at each grid point, values from neighbouring grid points
must be loaded multiple times (for different derivative stencil arms)
but written back only once per time step. The BSSN kernel reads each
of its 25 evolved fields from several neighbouring points for each
spatial derivative, leading to a high load-to-store ratio.

\subsection{Spill Fraction of Loads}

Within the ROI, the ratio of spills to total loads is:
\[
\frac{3{,}782{,}099{,}438}{12{,}477{,}164{,}927} \approx 30.3\%
\]
Nearly 1 in 3 load instructions is a spill reload. The BSSN evolution
kernel must simultaneously keep live a large number of floating-point
values at each grid point: the 25 evolved fields, their spatial
derivatives, intermediate expressions (Christoffel symbols, Ricci
tensor components, etc.), and loop variables. With RISC-V providing
32 floating-point registers, the compiler cannot avoid aggressive
spilling of these intermediates to the stack.

\subsection{ROI Coverage}

The ROI accounts for $48{,}023{,}254{,}961$ out of
$48{,}961{,}298{,}333$ total simulated instructions:
\[
\frac{48{,}023{,}254{,}961}{48{,}961{,}298{,}333} \approx 98.1\%
\]
This is the highest ROI coverage fraction observed across all
benchmarks. The instrumentation is near-perfect: \texttt{CCTK\_Evolve()}
encompasses essentially the entire meaningful computation of the
benchmark, with only 1.9\% of instructions spent on framework
initialisation and shutdown.

\subsection{Pre-ROI Spill Density}

The pre-ROI spill rate of 11.83\% ($110{,}850{,}678$ spills over
$936{,}822{,}146$ instructions) is itself the highest pre-ROI spill
rate observed across all benchmarks. This indicates that even the
Cactus framework's initialisation code --- which sets up the grid
structure, activates thorns, and computes initial data --- is
extremely register-pressure-intensive on RISC-V. The initial data
computation (GaugeWave setup) involves the same BSSN tensor
expressions as the evolution, explaining the similarly high spill density.

\subsection{Physical Interpretation of Register Pressure}

The BSSN equations in McLachlan form involve the simultaneous
computation of:
\begin{itemize}
    \item The conformal metric $\tilde{\gamma}_{ij}$ (6 independent
          components) and its 18 first spatial derivatives
    \item The conformal extrinsic curvature $\tilde{A}_{ij}$
          (6 independent components) and its derivatives
    \item The Christoffel symbols $\tilde{\Gamma}^k_{ij}$
          (18 independent components in 3D)
    \item The Ricci tensor $R_{ij}$ (6 independent components)
          and its contractions
    \item The gauge variables $W$, $K$, $\Gamma^i$, $\alpha$,
          $\beta^i$ and their derivatives
\end{itemize}
A single grid-point update in the ML\_BSSN kernel requires
$\mathcal{O}(100)$--$\mathcal{O}(1000)$ floating-point operations
with $\mathcal{O}(50)$ or more simultaneously live intermediate
values. The RISC-V F/D register file (32 floating-point registers)
is fundamentally insufficient to hold all live values simultaneously,
making spilling unavoidable regardless of compiler optimisation level.

\subsection{Scaling to Refrate}

The refrate dataset uses a $84^3$ grid (8$\times$ more grid points)
and 80 iterations (8$\times$ more time steps), making it $\sim$64$\times$
more expensive. Assuming spill rate remains constant, the estimated
refrate ROI spill count would be:
\[
3{,}782{,}099{,}438 \times 64 \approx 242 \text{ billion spills}
\]

% =========================================================
\section{Conclusion}

The \texttt{507.cactuBSSN\_r} benchmark exhibits the highest register
spill rate (7.88\%) among all SPEC CPU2017 benchmarks evaluated in
this study. With 3.78 billion spill events across 48 billion ROI
instructions, and 1 in 3 load instructions being a spill reload,
cactuBSSN places the most sustained register pressure of any benchmark
on the RISC-V floating-point register file. This is structurally
driven by the BSSN evolution kernel: evaluating the Einstein field
equations at each grid point requires $\mathcal{O}(50)$ simultaneously
live floating-point intermediates, far exceeding the 32-register
RISC-V F/D file capacity. The ROI captures 98.1\% of total execution,
making it an ideal representative of the full benchmark. These results
establish \texttt{507.cactuBSSN\_r} as the most register-spill-intensive
benchmark in the SPEC CPU2017 suite on RISC-V, and a critical workload
for evaluating any register file extension or compiler spill-reduction
strategy on this architecture.

\end{document}
